% Lettre de motivation (FR) — El Mehdi Haress
% Cible : Docteur en Simulation numérique & IA — Scalian Insights, Rennes (RedLUM / SIGMA)
% Compilation : xelatex lettre-scalian-simia.tex
\documentclass[10pt,a4paper]{article}

\usepackage{fontspec}
\usepackage[french]{babel}
\usepackage[T1]{fontenc}
\usepackage[
  top=1.5cm, bottom=1.3cm, left=2.0cm, right=2.0cm
]{geometry}
\usepackage{xcolor}
\usepackage{microtype}
\usepackage[hidelinks]{hyperref}
\usepackage{fontawesome5}

% ---------- Polices (identiques au CV) ----------
\setmainfont{texgyrepagella}[
  Extension      = .otf,
  Scale          = 1.05,
  UprightFont    = *-regular,
  BoldFont       = *-bold,
  ItalicFont     = *-italic,
  BoldItalicFont = *-bolditalic,
]
\newfontfamily\headingfont{texgyreheros}[
  Extension      = .otf,
  UprightFont    = *-regular,
  BoldFont       = *-bold,
  ItalicFont     = *-italic,
  BoldItalicFont = *-bolditalic,
]

% ---------- Couleurs (identiques au CV) ----------
\definecolor{accent}{RGB}{28,58,94}
\definecolor{ink}{RGB}{26,26,26}
\definecolor{muted}{RGB}{95,105,120}
\definecolor{rulegray}{RGB}{200,206,214}
\color{ink}

\newcommand{\cpp}{\mbox{C++}}

% ---------- Mise en page ----------
\pagestyle{empty}
\setlength{\parindent}{0pt}
\setlength{\parskip}{1.15ex plus 0.5ex minus 0.15ex}
\linespread{1.08}
\frenchspacing

\begin{document}

{\headingfont\bfseries\fontsize{20}{23}\selectfont EL MEHDI HARESS}\\[0.7ex]
{\headingfont\fontsize{10.5}{13}\selectfont\color{accent}%
Docteur en mathématiques appliquées \textbar{} Ingénieur CentraleSupélec}\\[1.0ex]
{\small
\faEnvelope[regular]~\href{mailto:el-mehdi.haress@inria.fr}{el-mehdi.haress@inria.fr}\quad
\faPhone*~+33 7 58 23 41 50\quad
\faLinkedin~\href{https://www.linkedin.com/in/el-mehdi-haress-67750b3aa/}{el-mehdi-haress}\quad
\faGlobe~\href{https://elmehdiharess.org}{elmehdiharess.org}
}

\vspace{0.4ex}
{\color{rulegray}\rule{\linewidth}{0.7pt}}

\vspace{1.1ex}

\begin{minipage}[t]{0.52\textwidth}
\vspace{0pt}\raggedright\small
Madame \textbf{Aurélie Collet}\\
Responsable Recrutement \& Relations Écoles\\
\textbf{Scalian} \textemdash{} Insights / Simulation \& IA\\
2 rue Antoine Becquerel, 35000 Rennes
\end{minipage}\hfill
\begin{minipage}[t]{0.42\textwidth}
\vspace{0pt}\raggedleft\small
Sophia Antipolis, le 16 septembre 2026
\end{minipage}

\vspace{1.8ex}

{\small\textbf{Objet :} candidature au poste de Docteur en Simulation numérique \& IA (H/F) \textemdash{} Rennes,
projets RedLUM et SIGMA}

\vspace{1.6ex}

Madame,

Ce qui m'a arrêté dans votre annonce n'est pas le titre, c'est le projet \textbf{RedLUM} : des
\textbf{modèles d'ordre réduit sous incertitude de localisation}, un couplage simulation--mesure, et une
paramétrisation stochastique des erreurs de réduction. J'ai passé cinq ans à démontrer des garanties
numériques pour des dynamiques stochastiques irrégulières. RedLUM est le terrain où ces garanties servent à
contrôler un système réel, avec des capteurs et une contrainte de temps réel.

Ma thèse construit des \textbf{schémas d'approximation pour des équations aux dérivées partielles
stochastiques à dérive singulière}, là où les hypothèses habituelles des schémas d'Euler tombent, avec
vitesses de convergence prouvées. À Leeds, j'ai ensuite quantifié la \textbf{sensibilité d'une solution aux
perturbations de sa spécification et de ses paramètres} : c'est, dans le langage de RedLUM, mesurer ce que
coûte une réduction, et sous quelles hypothèses elle se dégrade. Je suis aujourd'hui ingénieur de recherche à
l'\textbf{Inria}, équipe CALISTO, sur des méthodes numériques probabilistes pour des systèmes stochastiques
multi-échelles issus d'applications industrielles. Collaborer avec Inria et l'INRAE sur un ANR, tout en
portant le volet logiciel chez Scalian, est exactement le format de travail que je cherche.

Le second projet, \textbf{SIGMA}, m'intéresse par l'autre bout : apprendre une loi de commande dans un
environnement simulé, puis la déployer. Je ne prétends pas arriver avec Isaac Sim ou de l'apprentissage par
renforcement déjà en production. En revanche, j'ai publié à \emph{AISTATS} sur l'apprentissage sous
\textbf{contrainte de risque} (mesures spectrales, queues lourdes), je code en \textbf{Python} au quotidien
depuis huit ans, et j'ai les bases du \cpp{} et de MATLAB de ma formation d'ingénieur. Monter en compétence
sur ITHACA-FV, OpenFOAM, le MLOps et le déploiement embarqué est un travail que je veux faire, pas une
ligne que j'écris pour plaire.

Enfin, votre offre demande de \textbf{publier, présenter, et former}. J'ai huit publications internationales,
plus de quinze exposés, enseigné la finance stochastique et les EDP à CentraleSupélec, et co-organisé deux
rencontres scientifiques. Je suis mobile sur Rennes à compter de novembre.

Je serais très heureux d'en parler avec Valentin et Florian : ce que la stochastique des modèles réduits
peut apporter à RedLUM, et ce qu'il me reste à apprendre pour SIGMA.

Veuillez agréer, Madame, l'expression de mes salutations distinguées.

\vspace{2.2ex}

\hfill
\begin{minipage}{0.42\textwidth}
\raggedleft
{\headingfont\bfseries\small El Mehdi Haress}\\[0.2ex]
{\small\color{muted}Docteur en mathématiques appliquées}
\end{minipage}

\end{document}
